\section*{Bab \ref{ch:b1:organomagnesium} --- Pereaksi Organomagnesium}

\begin{solution}{exo:b1:organomagnesium:1}
\ce{CH3CH2Br + Mg -> CH3CH2MgBr}; \ce{C6H5Br + Mg -> C6H5MgBr} (dalam eter
kering). C--Br: Br (2.96) lebih elektronegatif daripada C (2.55), karbon
positif. C--Mg: Mg (1.31) kurang elektronegatif, karbon negatif.
\end{solution}

\begin{solution}{exo:b1:organomagnesium:2}
Dengan air: \ce{CH3MgI + H2O -> CH4 + Mg(OH)I}.
Dengan etanol: \ce{CH3MgI + CH3CH2OH -> CH4 + CH3CH2OMgI}.
Dengan asam etanoat: \ce{CH3MgI + CH3COOH -> CH4 + CH3COOMgI}.
Dengan deuterium oksida: \ce{CH3MgI + D2O -> CH3D + Mg(OD)I}.
\end{solution}

\begin{solution}{exo:b1:organomagnesium:3}
Metanal: propan-1-ol (primer). Etanal: butan-2-ol (sekunder). Propanon:
2-metilbutan-2-ol (tersier). Benzaldehida: 1-fenilpropan-1-ol (sekunder).
Karbon dioksida: asam propanoat.
\end{solution}

\begin{solution}{exo:b1:organomagnesium:4}
Gugus OH-nya cukup asam untuk merusak pereaksi segera setelah pereaksi
itu terbentuk: molekul itu akan memprotonasi dirinya sendiri (atau
tetangga-tetangganya), menghasilkan magnesium alkoksida butan-1-ol,
\ce{CH3(CH2)3OMgBr}, bukan pereaksi yang dapat dipakai. OH itu harus
dilindungi lebih dahulu (\cref{ch:b1:acetals-protection} memperlihatkan
gagasannya untuk aldehida).
\end{solution}

\begin{solution}{exo:b1:organomagnesium:5}
Pasangan C--Mg menyerang karbon karbonil etanal, dengan pasangan $\pi$ C=O
berpindah ke O: \ce{CH3CH(CH3)O^-} dengan \ce{MgBr+}; lalu \ce{H3O+}
memberikan sebuah proton kepada oksigen: propan-2-ol. Propan-2-ol membawa
dua gugus metil pada karbon karbinolnya: zat itu \omterm{def:b1:stereochemistry-in-depth:stereocentre}{akiral}. (Dengan R yang
lain, misalnya etilmagnesium bromida pada etanal, produknya, butan-2-ol,
\omterm{def:b1:stereochemistry-in-depth:stereocentre}{kiral} tetapi rasemat: karbonil yang datar diserang pada kedua sisinya
sama banyak.)
\end{solution}

\begin{solution}{exo:b1:organomagnesium:6}
Karbon karbinol membawa satu gugus metil dan dua gugus etil:
etilmagnesium bromida dengan butan-2-on, atau metilmagnesium bromida
dengan pentan-3-on.
\end{solution}

\begin{solution}{exo:b1:organomagnesium:7}
Buatlah \ce{(CH3)2CHMgBr} dari 2-bromopropana dan magnesium dalam eter
kering, tuangkan ke atas \ce{CO2} padat, lalu asamkan: \ce{(CH3)2CHCOOH}.
$n = 12.3/123.0 = 0.100$ mol;
$0.100 \times 0.70 \times 88.0 = \qty{6.2}{g}$.
\end{solution}

\begin{solution}{exo:b1:organomagnesium:8}
$0.18/18.0 = \qty{0.010}{mol}$ air merusak \qty{0.010}{mol} pereaksi:
\qty{20}{\percent}; terbentuk benzena, \ce{C6H5MgBr + H2O -> C6H6 +
Mg(OH)Br}.
\end{solution}

\begin{solution}{exo:b1:organomagnesium:9}
Pereaksi itu, sebuah \omterm{def:b1:electronic-effects:nucleophile}{nukleofil}, bereaksi dengan bromobenzena yang masih
ada: secara keseluruhan \ce{C6H5MgBr + C6H5Br -> C6H5-C6H5 + MgBr2}.
Menambahkan bromida perlahan-lahan ke dalam magnesium berlebih menjaga
konsentrasinya tetap rendah: bromida itu bertemu logam yang segar,
bukan pereaksi yang sudah terbentuk.
\end{solution}

\begin{solution}{exo:b1:organomagnesium:10}
1-Feniletan-1-ol: \ce{CH3MgBr} (dari bromometana) dengan benzaldehida,
atau \ce{C6H5MgBr} (dari bromobenzena) dengan etanal.
2-Fenilpropan-2-ol: \ce{C6H5MgBr} dengan propanon, atau \ce{CH3MgBr}
dengan 1-feniletan-1-on.
\end{solution}

\begin{solution}{exo:b1:organomagnesium:11}
Keton ditambahkan ke dalam pereaksi agar pereaksi, yang dibuat di dalam
labu yang kering, tidak pernah meninggalkannya dan kalor yang dilepaskan
oleh setiap porsi diserap oleh eter yang mendidih. Air saja akan
menghasilkan garam basa \ce{Mg(OH)Br} yang bergelatin, yang memerangkap
produknya; asam encer yang dingin melarutkannya sebagai \ce{Mg^2+}, dan
dinginnya membatasi reaksi samping alkohol tersier dalam asam.
\end{solution}

\begin{solution}{exo:b1:organomagnesium:12}
Setiap mol pereaksi menghasilkan dengan air satu mol garam magnesium
basa, \ce{RMgBr + H2O -> RH + Mg(OH)Br}, yang \ce{OH-}-nya dititrasi oleh
asam: jumlah pereaksi sama dengan asam yang terpakai, \qty{0.088}{mol},
\omterm{def:b1:extent-q-and-k:fractional-extent}{rendemen} \qty{88}{\percent}.
\end{solution}

\begin{solution}{pb:b1:organomagnesium:1}
\textbf{1.} \ce{C6H5Br + Mg -> C6H5MgBr}.
\textbf{2.} Mg $0.535/24.3 = \qty{0.0220}{mol}$; \ce{C6H5Br} $3.14/156.9 =
\qty{0.0200}{mol}$: magnesium sedikit berlebih.
\textbf{3.} Air merusak pereaksi: \ce{C6H5MgBr + H2O -> C6H6 +
Mg(OH)Br}.
\textbf{4.} Tabung itu mencegah uap air dari udara masuk melalui
kondensor.
\textbf{5.} Eter memberikan \omterm{def:b1:lewis-resonance-vsepr:lewis}{pasangan elektron bebas} kepada magnesium
(\omterm{def:b1:lewis-resonance-vsepr:lewis-acid}{basa Lewis}): tanpa eter pereaksi itu tidak terbentuk atau tidak tetap
larut.
\textbf{6.} \ce{C6H5MgBr + C6H5Br -> C6H5-C6H5 + MgBr2}; bifenil
$0.15/154.0 = \qty{9.7e-4}{mol}$, yang memakai \qty{9.7e-4}{mol}
bromobenzena dan pereaksi dalam jumlah yang sama: seluruhnya
\qty{1.9e-3}{mol} bromobenzena hilang.
\textbf{7.} $3.28/182.0 = \qty{0.0180}{mol}$.
\textbf{8.} Pasangan C--Mg menyerang karbon karbonil, pasangan $\pi$
berpindah ke oksigen: \ce{(C6H5)3C-O^-} \ce{MgBr+}.
\textbf{9.} Pereaksi yang tersedia paling banyak
$0.0200 - 0.0019 = \qty{0.0181}{mol}$; benzofenon \qty{0.0180}{mol}:
benzofenon yang membatasi, dengan selisih tipis.
\textbf{10.} Tidak: karbon karbinol membawa tiga gugus fenil yang
identik.
\textbf{11.} Pereaksi tetap di dalam labunya yang kering, dan kalor setiap
porsi keton diserap oleh eter yang mendidih.
\textbf{12.} Air itu akan merusak sebagian pereaksi (benzena) dan
menyisakan benzofenon yang tidak bereaksi.
\textbf{13.} \ce{(C6H5)3COMgBr + H3O+ -> (C6H5)3COH + Mg^2+ + Br- + H2O}.
\textbf{14.} Asam melarutkan garam-garam magnesium basa yang akan
diendapkan oleh air saja; dinginnya membatasi pembentukan \omterm{def:b1:electronic-effects:intermediates}{karbokation}
pada pertanyaan 18.
\textbf{15.} Trifenilmetanol dan bifenil di lapisan eter; garam-garam
magnesium di lapisan air.
\textbf{16.} Mencuci (atau menggerus) padatan itu dengan petroleum eter
yang dingin: bifenil larut, trifenilmetanol tertinggal.
\textbf{17.} Titik leleh yang tajam dan sama dengan nilai tabel (atau satu
noda pada kromatografi lapis tipis).
\textbf{18.} \ce{(C6H5)3COH + H+ -> (C6H5)3C+ + H2O}: \omterm{def:b1:electronic-effects:intermediates}{karbokation} tersier
yang muatannya tersebar ke tiga cincin benzena melalui resonansi, sangat
stabil, dan berwarna karena konjugasi yang luas itu.
\textbf{19.} $3.75/260.0 = \qty{0.0144}{mol}$.
\textbf{20.} $0.0144/0.0180 = \qty{80}{\percent}$.
\textbf{21.} Kehilangan pada pemindahan dan rekristalisasi; pereaksi yang
rusak oleh jejak kelembapan; reaksi yang tidak sempurna.
\textbf{22.} Asam benzoat, \ce{C6H5COOH}, sesudah pengasaman.
\textbf{23.} $0.0181 \times 122.0 = \qty{2.21}{g}$.
\textbf{24.} \ce{C6H5MgBr + CO2 -> C6H5COOMgBr}, lalu
\ce{C6H5COOMgBr + H3O+ -> C6H5COOH + Mg^2+ + Br- + H2O}.
\textbf{25.} \textbf{\qty{80}{\percent}}.
\end{solution}
